\section{High-frequency trading: strategies, speed races and the costs of other traders}

\begin{qfdefinition}[High-frequency trading: proprietary strategies, the speed tools they share, and why the label names no single mechanism]\label{def:hft-strategies}
\qffield{Statement} In the SEC's 2010 concept release the term ``high frequency trading'' is ``relatively new and is not yet clearly defined''. It typically refers to professional traders acting in a proprietary capacity whose strategies generate many trades a day, organised as proprietary trading firms, proprietary desks of broker-dealers or hedge funds \citep[sec.~IV.B, p.~45]{sec-concept-release-2010}. Five characteristics are often attributed to them: (1) high-speed programs for generating, routing and executing orders; (2) co-location and individual data feeds to cut latency; (3) very short holding periods; (4) many orders cancelled shortly after submission; (5) ending the day close to flat (p.~45; restated in \citep[sec.~2.4.3, p.~163]{alpha-laruelle}). Rather than define HFT, the release treats \emph{strategies} and \emph{tools} separately. One firm may run many strategies at once, some generating many trades and some not, and ``conceivably'' some benefit market quality and long-term investors while others harm them \citep[sec.~IV.B, p.~46]{sec-concept-release-2010}. To present issues for comment, ``without limiting the broad request'', it describes four broad types of strategy often associated with proprietary firms \citep[sec.~IV.B.1, pp.~48--57]{sec-concept-release-2010}:
\begin{itemize}
\item \emph{passive market making}: mainly non-marketable resting bids and offers, profiting from the spread and from liquidity rebates, and taking liquidity only to liquidate a position quickly. Layering many orders can give cancellation rates of 90 percent or more, with orders often lasting a second or less (pp.~48--49);
\item \emph{arbitrage}: capturing price discrepancies between related products or markets, such as an ETF against its basket, mostly by taking liquidity, with hedged positions that may last several days. The release asks how far such strategies trade price differences between trading centres and depend on latency between direct and consolidated data feeds (pp.~51--52);
\item \emph{structural}: exploiting structural vulnerabilities in the market or in other participants. Examples are using the fastest data to trade with participants still offering executions at stale prices, and setting a new NBBO with a small order and then triggering guaranteed-match trades on the other side (pp.~52--53);
\item \emph{directional}: a significant, unhedged position on an expected intraday move. Some, such as trading a temporary departure from fundamental value, ``often may contribute'' to price discovery. The release singles out two: \emph{order anticipation}, detecting large buyers or sellers and trading ahead of them, and \emph{momentum ignition}, orders and trades, perhaps with false rumours, meant to spark a rapid price move or trigger stop orders, followed by unwinding the position (pp.~53--57).
\end{itemize}
Many of these strategies are not new; what is new is the technology for identifying and executing them (p.~48). The tools discussed separately are co-location and individual trading-centre data feeds. Speed matters ``both in the absolute sense of achieving very small latencies and in the relative sense of being faster than competitors'' (sec.~IV.B.2, pp.~57--58).

\qffield{Conditions} U.S. equities as of January 2010; European venues as of the 2018 printing; Narang's account as of 2024; FINRA Rule~5270 as displayed on 15 October 2026, the only current rule checked; no market share checked and the historical figures dated as given. Interpretation: speed and automation are means shared by these strategies, not a strategy, and the economic mechanism differs between them. Passive market making is paid the spread for supplying immediacy and bears inventory and adverse-selection risk (\hyperref[mod:avellaneda-stoikov]{Avellaneda--Stoikov}, \hyperref[meth:dealer-quote-adjustment]{dealer quote adjustment}). Arbitrage enforces relations between prices (\hyperref[def:arbitrage]{arbitrage}). Trading against stale quotes earns rents from reacting first to public information (\hyperref[res:latency-arbitrage-arms-race]{latency-arbitrage arms race}). Order anticipation profits from information about others' orders rather than about values (\hyperref[def:order-anticipation]{order anticipation}). Momentum ignition is a form of bluffing (\hyperref[def:bluffing-manipulation]{bluffing and manipulation}). Effects on liquidity, price discovery and other traders' costs must therefore be judged strategy by strategy.

The release's four types organise a request for comment on U.S. equities in January 2010; they are not a universal classification, and they overlap: trading against stale prices is filed under ``structural'' and cross-venue price differences under ``arbitrage''.

\qffield{Practice} Lehalle and Laruelle describe high-frequency market makers in European equities around 2009--2012 as arbitraging between the venues that quote the same stock, which keeps those prices consistent. Maker--taker fee schedules on MTFs reward this, and trading very small orders very often lowers their inventory risk \citep[secs.~2.4 and~2.4.3, pp.~158 and~163]{alpha-laruelle}. Venues court them with low latency and co-location, ``especially since the lower the latency, the smaller the inventory risk''; the authors read one venue's low ``flow reactivity'' as its high-frequency market makers offering less liquidity when activity rises (p.~158). Probing orders are a tool, not a strategy. The release defines a ``pinging'' order as an immediate-or-cancel order used to find undisplayed liquidity, and distinguishes its use in a normal search for liquidity from its use to detect and trade in front of large interest \citep[p.~47, n.~69]{sec-concept-release-2010}.

An order anticipator that detects a large buyer first buys, and is largely indifferent to whether the large trader is buying or selling; a long-term investor seeking liquidity wants the opposite side, and both parties benefit from that search \citep[p.~55]{sec-concept-release-2010}. Liquidity seeking by routers is in \hyperref[def:smart-order-routing]{smart order routing}. Narang groups high-frequency strategies as contractual and non-contractual market making, arbitrage and \emph{fast alpha}. Fast alpha applies the price-based alphas of slower strategies, momentum, mean reversion and technical sentiment, to prices, volumes and order-book information that change during the day; its statistical relationships carry risk because ``there is no structural reason'' they must persist, and intraday trend following is less common than reversion because impact and slippage work against it \citep[chap.~15, ``High-Frequency Trading'', sec.~15.4, pp.~291--293]{s7-narang}.

Execution algorithms also adjust their aggressiveness to the order book's imbalance, a ``micro price'' (chap.~7, pp.~133--134), and one that leaves a detectable footprint lets other participants react ``to increase the market impact and slippage'' it incurs (p.~130). \emph{Fragmentation and roles}. Lehalle and Laruelle call HFT ``the price to pay for fragmentation; it is not possible to put trading venues in competition without agents building high-frequency liquidity bridges across them'', and describe high-frequency traders as ``acting as market-makers but in most cases without obligations'', which blurred the roles of each layer of the market \citep[Introduction, pp.~5--6]{alpha-laruelle}; HFT updates books faster than a non-co-located observer's round trip, so an order's role as provider or taker is known only in probability (p.~22).

The identification criteria (automated generation, routing and cancellation; low latency, co-location and direct feeds; short holding periods and low overnight inventory; high message-to-trade ratios) restate the release's characteristics above; that they ``describe technology and horizon, not whether the strategy is beneficial'' is the interpretation's caveat, consistent with the release.

\qffield{Pitfalls} A proposed taxonomy lists five strategies: passive market making, statistical arbitrage, latency arbitrage, order anticipation and liquidity detection. Compared with the release:
\begin{itemize}
\item ``statistical arbitrage'' on short-horizon relative value or lead--lag signals straddles the arbitrage and directional types;
\item ``latency arbitrage'' corresponds to the structural example of trading against stale prices, and across venues to the release's arbitrage question;
\item ``liquidity detection'' is a tool that order anticipators and ordinary liquidity seekers both use, not a strategy;
\item momentum ignition, which the release names, was missing.
\end{itemize}
The statement, that one firm may run several strategies whose risk and social value differ, matches p.~46, where the difference in social value is a possibility to assess, not a finding. A strategy label does not establish illegality. The release treats order anticipation that involves no breach of duty or misappropriation as a question of policy; front running in breach of a duty, false rumours and manipulative conduct were already prohibited (pp.~53--55 and~57). These legal remarks are as of January 2010 and are not statements of current law. The release's estimate that HFT was 50 percent or more of U.S. equity volume came from press reports (p.~45, n.~67).

Lehalle and Laruelle's ``more than 80\% of the deals in the US, 40\% in Europe, and 35\% in Japan'' is reported as hearsay (``They are said to be'', p.~163) without source or date. Neither is a measured share. \emph{Trade counts as a proxy}. Lehalle and Laruelle read a rising number of trades with a falling average trade size as evidence of growing high-frequency activity: ``the number of trades per day is a good proxy to get an idea of how intense high frequency activity is on a stock'' (FTSE~100 and CAC~40 rising, IBEX~35 flat, 2007--2012) \citep[sec.~2.5.1, pp.~170--171, Figure~2.25; introduction, pp.~20--21, Figure~5]{alpha-laruelle}.

It is a proxy, not a measure: the caveat that the same count can reflect genuine liquidity, order slicing or stress-induced churn is consistent with the 2010 staff report's warning that ``high trading volume is not necessarily a reliable indicator of market liquidity'' (\hyperref[ev:may-2010-flash-crash]{the May 2010 Flash Crash}), and no venue's HFT share was checked. ``Microstructure alphas'' on bets of seconds or less correspond to Narang's fast alpha; its claim that machine learning can discern other players' execution patterns, more easily the worse their execution models, and that such patterns are ``somewhat stable'', is not in the third edition and has no support in the passages read.

Its ``HFT shark strategy'' (sending very small orders to detect a large iceberged order, then trading in front of it and exiting ``with a quick and relatively riskless profit'') is also not in the third edition; it describes pinging used for order anticipation, as the release does (p.~47, n.~69; p.~55). Three qualifications. (i) ``Front-run'' is used loosely: Narang defines front-running as a fiduciary using knowledge of a customer's order before acting on it, and argues that most high-frequency traders, who see orders only once they reach the market, do not front-run in that sense, except contractual market makers, who see customer flow and for whom it is ``theoretically possible'' (p.~312) \citep[Appendix, sec.~A.2.1, pp.~312--314]{s7-narang}; FINRA's rule bars members from trading with ``material, non-public market information concerning an imminent block transaction'' \citep[para.~(a)]{finra-rule-5270}, and inferring a hidden order from fills is not obviously that information on the rule's wording (Interpretation; the term is not defined there, FINRA guidance was not checked, and the rule's supplementary material notes that other rules, such as Rules~2010 and~5320, may apply to front-running other orders).

Whether pinging for order anticipation is lawful elsewhere was not checked. (ii) The profit is not riskless: the anticipator bears the risk that the inferred order is not there, is smaller or ends, and that the price does not move (Interpretation; \hyperref[def:order-anticipation]{order anticipation}). (iii) The block is a historical description, kept with its original wording in the artefact, not a supported mechanism.

\qffield{See also} \hyperref[res:latency-arbitrage-arms-race]{latency-arbitrage arms race}; \hyperref[res:hft-spread-investor-cost]{spread against investor cost}; \hyperref[def:order-anticipation]{order anticipation} and \hyperref[def:bluffing-manipulation]{bluffing and manipulation}; \hyperref[def:arbitrage]{arbitrage}; \hyperref[def:smart-order-routing]{smart order routing} (owned by \nolinkurl{market-structure-execution-quality}).

\qffield{Further reading} Menkveld's survey (\nolinkurl{wide-menkveld-hft-review}), the pathway's first reading, was not available and is unread.

\end{qfdefinition}

\begin{qfresult}[Latency arbitrage in a continuous limit order book: stale quotes are sniped, the speed race dissipates the rents, and investors pay through the spread (Budish, Cramton and Shim)]\label{res:latency-arbitrage-arms-race}
\qffield{Statement} Budish, Cramton and Shim's model \citep[secs.~VI.A--VI.B, pp.~1576--1583]{wide-budish-cramton-shim}. A security $x$ trades on a continuous limit order book. Its value is perfectly correlated with a public signal $y$, which jumps as a compound Poisson process with rate $\lambda_{\mathrm{jump}}$ and absolute jump size $J$, and $x$ can always be liquidated at that value. Investors arrive at rate $\lambda_{\mathrm{invest}}$, each needing to buy or sell one unit, and take liquidity. Risk-neutral trading firms make and take liquidity, and messages that arrive together are processed serially in random order. In equilibrium exactly one unit is quoted on each side, a bid at $y-s/2$ and an ask at $y+s/2$; the two quotes may belong to one trading firm or to two different firms (Proposition~1(i), p.~1583).

The authors describe play with one firm in the liquidity-provider role for simplicity, noting economically equivalent equilibria in which the role rotates or is random (pp.~1579--1580). When $y$ jumps by more than $s/2$, the firm whose quote is stale sends a cancel while the other $N-1$ firms send orders to trade against (``snipe'') it, so the stale quote is sniped with probability $(N-1)/N$ (Proposition~1(iii), p.~1583). The spread equates investors' payments with the rents from sniping:
\[ \lambda_{\mathrm{invest}}\,\frac{s}{2}=\lambda_{\mathrm{jump}}\Pr\Bigl(J>\frac s2\Bigr)\,\mathbb E\Bigl[J-\frac s2\Bigm|J>\frac s2\Bigr] \]
(eq.~(3), p.~1582), which does not depend on $N$. The spread is positive with symmetric information and no inventory costs. Sniping costs grow linearly with the quantity quoted, while the benefit of quoting more grows less than linearly, so the spread widens unit by unit and the book is thin (Proposition~2, pp.~1583--1585). This assumes investors trade their whole quantity at once, which the authors note is not optimal for them. When firms can rent speed at $c_{\mathrm{speed}}$ per unit time, the zero-profit conditions for the liquidity provider and the snipers (eqs.~(5)--(6), p.~1589) give
\[ \lambda_{\mathrm{invest}}\,\frac{s^*}{2}=N^*c_{\mathrm{speed}} \]
(eq.~(7), p.~1590): the total prize, investors' cost through the spread and the total spent on speed are equal, so the rents are dissipated by the race (Proposition~3, pp.~1590--1591; page images of pp.~1589--1590). The race is a prisoner's dilemma: welfare would be higher by $N^*c_{\mathrm{speed}}$ if the firms could commit not to invest (Proposition~4, p.~1591). The prize rises with the frequency of jumps and with a mean-preserving spread of their distribution, and does not depend on the cost of speed, the speed of the fastest firms or the speed difference (Proposition~5, p.~1592).

\qffield{Notation} $\lambda_{\mathrm{jump}},\lambda_{\mathrm{invest}}$ are Poisson arrival rates, not impact slopes; $s$ is the full bid--ask spread in price units; $N$ is the number of fast trading firms; $c_{\mathrm{speed}}$ is a rental cost per unit time. Starred values are equilibrium values with endogenous entry.

\qffield{Conditions} The assumptions are deliberately strong, a ``best case'' for the continuous book: a perfectly correlated public signal, costless liquidation, no asymmetric information or inventory costs, unit investor demand, prices on a continuum, and equally fast fast firms (pp.~1576--1579). With a discrete tick the liquidity-provider role is strictly preferred and firms race for the top of the book instead (n.~19, p.~1580). The authors do not claim that markets are less liquid than before electronic trading and HFT, which they say the evidence contradicts. Their claim is that markets are less liquid than under a design without sniping, the frequent batch auctions they propose (n.~5, pp.~1554--1555; abstract, p.~1548). Evidence (sec.~I, pp.~1552--1553, summarising sec.~V, which was not read): ES futures and SPY, 2005--2011, exchange direct-feed data.

In 2011 the median return correlation was 0.1016 at 10~ms and 0.0080 at 1~ms. The median duration of ES--SPY arbitrage opportunities fell from 97~ms in 2005 to 7~ms in 2011, while the median profit, about 0.08 index points per unit traded, was roughly constant, and their frequency moved with volatility. The ES--SPY prize averaged \$75 million a year. The authors call extrapolation to other instruments ``common sense'' and decline to give a total.

\qffield{Practice} With speed costs that differ across firms, the marginal fast firm earns zero profit and inframarginal firms earn strictly positive profits; total sniping rents then exceed total spending on speed (n.~27, p.~1590). Lehalle and Laruelle describe a related competitive margin. High-frequency market makers extended coverage to less liquid stocks once their hedging was built and their technology, a fixed cost, was paid; growing competition between them ``forced them to diversify their sources of revenue''. Below a liquidity threshold, hedging and short-selling costs keep them out \citep[sec.~2.4.2, pp.~161--162]{alpha-laruelle}. This is their interpretation of coverage statistics for European indexes, 2009--2012, not a test.

\qffield{Pitfalls} The assertion was that that speed advantages are valuable only while competitors cannot replicate them, and that once speed diffuses, durable profits must come from inventory control, better models, lower costs or privileged order flow. In this model and its ES--SPY evidence, competition does not shrink the arbitrage prize; it raises the speed needed to win it, and under free entry the rents go into speed spending, which investors bear through the spread. What erodes is a firm's \emph{relative} lead: at the margin, profit net of speed costs is zero (Interpretation of eqs.~(5)--(7), checked algebraically and numerically). Lower costs of speed do sustain profits (n.~27). Inventory control, better models and privileged order flow are not in the passages read and remain unsupported assertions, kept in the artefact. The empirical constancy is for one pair, 2005--2011. The pathway's evidence on races at a UK venue (\nolinkurl{exp-aquilina-budish-oneill-races}) and on HFT profits and competition (\nolinkurl{wide-baron-hft-competition}) is unread.

\qffield{See also} \hyperref[def:hft-strategies]{HFT strategies}; \hyperref[res:hft-spread-investor-cost]{spread against investor cost}; \hyperref[mod:glosten-milgrom]{Glosten--Milgrom} (adverse selection from private information, against sniping on public information); \hyperref[def:arbitrage]{arbitrage}; \hyperref[def:smart-order-routing]{smart order routing} (latency and mirrored liquidity).

\qffield{Further reading} The frequent-batch-auction analysis is in secs.~VII--VIII of the paper.

\end{qfresult}

\begin{qfresult}[A narrower spread need not mean lower costs: Lehalle and Laruelle's stylised model of a fast trader quoting inside the book]\label{res:hft-spread-investor-cost}
\qffield{Statement} Lehalle and Laruelle's stylised book \citep[sec.~2.4.3, pp.~164--165, Figure~2.21, page images]{alpha-laruelle}. Two traders, A and B, rest at the best bid and ask, a spread $\psi$ apart, and take turns becoming impatient. Each pays $\psi$ when crossing and nothing when the other crosses, so each pays $\psi/2$ per trade on average. A high-frequency trader then quotes just inside B, so that A, crossing, pays $\psi/2+\epsilon$; it at once steps in front of A, and B pays $\psi/2+\epsilon$ too. The effective spread falls to a little over $\psi/2$, but A and B now pay $\psi/2+\epsilon$ on every trade, because each must cross every time, while the fast trader pays nothing. In the authors' words, ``the effective spread is reduced on an overall basis, but \dots\ costs for other participants can remain unchanged or slightly increased'' (p.~164).

In the calibrated version, $\psi_B$ is the spread before HFT and $\psi_{\mathrm{HFT}}$ the tighter one an HFT quotes. An investor hitting an HFT pays $\psi_B/2+\epsilon$ with $\epsilon=\psi_{\mathrm{HFT}}/2$, and if HFTs are in a fraction $q$ of trades the spread becomes $\psi_N=q\psi_{\mathrm{HFT}}+(1-q)\psi_B$ (p.~165). An investor posting one order meets an HFT with probability $r$ and then pays $(\psi_{\mathrm{HFT}}+\psi_B)/2$. Otherwise, with equal probability, the investor crosses and pays $\psi_B$ or stays passive and pays 0, so
\[ \mathbb E[C]=\frac{\psi_B+r\,\psi_{\mathrm{HFT}}}{2},\qquad \operatorname{Var}(C)=r\Bigl(\frac{\psi_{\mathrm{HFT}}+\psi_B}{2}\Bigr)^2+(1-r)\frac{\psi_B^2}{2}-\mathbb E[C]^2 \]
(Figure~2.22, p.~166, page image). The average cost of $N$ independent trades has the same mean and variance $\operatorname{Var}(C)/N$ (p.~168). Meeting HFTs more often raises the expected cost and lowers its dispersion. The dispersion falls a lot for an investor with few trades, such as a retail investor, and little for one with many, such as a pension fund (pp.~167--169, Figures~2.23--2.24).

\qffield{Notation} $\psi$ is a spread in basis points of price; $C$ is one order's cost. Interpretation: in the figure each trader's cost is measured from its own resting quote, so a passive fill costs 0 and a crossing the full spread. $C$ is therefore a cost against waiting at one's own quote, not the half-spread against the midpoint of \hyperref[def:spread-measures]{spread measures}.

\qffield{Conditions} Calibration: $\psi_B=5.41$~bp (April to November 2007) and $\psi_N=4.73$~bp (April to November 2011), for CAC~40 stocks on the primary market, Euronext Paris, only, with $q=40$ percent assumed (pp.~165--166, page image of p.~166). Interpretation, recomputed with exact arithmetic: these imply $\psi_{\mathrm{HFT}}=(\psi_N-(1-q)\psi_B)/q=3.71$~bp. $\mathbb E[C]$ then runs from 2.705~bp at $r=0$ to 4.56~bp at $r=1$, which matches the ends of Figure~2.24, and is 3.26~bp at $r=30$ percent and 3.63~bp at $r=50$ percent. The model attributes the whole fall in the primary-market spread between 2007 and 2011 to HFT and holds $\psi_B$ fixed, as if the earlier spread would otherwise have persisted.

It fixes $q$ rather than measuring it, ignores trading moved to other venues, and assumes that an investor who meets an HFT must cross while others cross half the time. The two spreads are observations; their attribution to HFT and the cost model are the authors' stylised interpretation, not an estimate of HFT's effect.

\qffield{Practice} The caveat's advice is consistent with this: judge fast liquidity provision by what other traders pay (effective and realised spreads, \hyperref[def:spread-measures]{spread measures}), by depth, adverse selection and behaviour under stress, not by quoted spreads or message counts. The SEC's 2010 release asked these as open questions. Are orders lasting a second or less ``phantom liquidity that disappears when most needed'', or do many firms quoting together give a stable market with many size updates but few changes in the NBBO price? Do proprietary firms, which have replaced specialists and market makers with affirmative and negative obligations, supply liquidity in times of stress \citep[sec.~IV.B.1, pp.~48--51]{sec-concept-release-2010}? In Budish, Cramton and Shim's model the cost of being sniped is priced into the spread and grows with the depth shown (\hyperref[res:latency-arbitrage-arms-race]{latency-arbitrage arms race}).

Lehalle and Laruelle add that HFTs' preference for small trades makes other traders split orders and bear more reporting and technology costs, ``fixed losses'' unevenly distributed (p.~165); this is an argument, not a measurement.

\qffield{Pitfalls} The text says that raising the probability of meeting an HFT from 30 to 50 percent ``increases the expected cost by approximately 0.5bp'' (p.~168). The book's own formula and calibration give $0.2\times3.71/2=0.37$~bp, and Figure~2.24's markers sit at about 3.27 and 3.64~bp, which agrees with 0.37 (page image; Interpretation). Its ``by 20\%'' means 20 percentage points. The caveat's ``high message rates can improve displayed prices'' and ``can also increase cancellation risk and stale-book uncertainty'' are, in the passages read, a stylised mechanism and open questions, not measured effects. The pathway's empirical readings on liquidity and price discovery (\nolinkurl{wide-hendershott-jones-menkveld}, \nolinkurl{wide-brogaard-hft-discovery}) are unread. The model's conclusion concerns expected cost per trade under its assumptions; it is not a finding that HFT raised investors' costs.

The caveats ``tighter spread, different cost'' and ``investor size asymmetry'' (a small investor bears high variance over a few trades, a large one diversifies it across child orders but still pays the aggregate cost), restate this entry's result; the size asymmetry is an interpretation and is not in the pages checked.

\qffield{See also} \hyperref[def:hft-strategies]{HFT strategies}; \hyperref[res:latency-arbitrage-arms-race]{latency-arbitrage arms race}; \hyperref[def:spread-measures]{spread measures} and \hyperref[res:spread-determinants]{spread determinants}; \hyperref[def:smart-order-routing]{smart order routing} (displayed liquidity is only probable); \hyperref[def:order-anticipation]{order anticipation} (quote matching steps ahead of passive traders).

\end{qfresult}

\begin{qfresult}[The value of queue position: speed buys priority, and in large-tick books a place at the front is worth a sizeable fraction of the spread (Moallemi and Yuan)]\label{res:queue-position-value}
\qffield{Statement} Under price--time priority, a resting order's value depends on its place in the first-in-first-out queue at its price. Moallemi and Yuan's model \citep[author working paper, June 2017 revision, abstract and sec.~1, pp.~1--4]{wide-moallemi-yuan-queue} identifies two components:
\begin{itemize}\setlength{\itemsep}{0pt}
\item a \emph{static} component, the trade-off at execution between earning the spread and bearing adverse selection, ``and incorporates the fact that adverse selection costs are increasing with queue position'': an order at the back of a long queue executes in the next instance only against a large trade, and large trades come more often from informed traders, so ``a good queue position acts as a filter on the population of contra-side market orders'' (after Glosten 1994; pp.~1 and~3);
\item a \emph{dynamic} component, ``the optionality associated with the future value that accrues by locking in a given queue position'', as the order moves up after trades and cancellations (pp.~1 and~4).
\end{itemize}
Speed is how positions are won: ``orders that are emitted faster will wind up in a better queue position, and this is especially important when traders are reacting to common signals'', notably right after a price change, when there is ``a race to establish queue positions at the new price'' (p.~3).

\qffield{Conditions} The model is for \emph{large-tick} assets, those for which ``the bid-ask spread rarely exceeds the minimum tick size'' (after Eisler et al., 2012): moving the price by a tick is costly, so traders queue rather than improve the price, and queues are long; for small-tick assets priority can be bought with a better price and queue position matters less (sec.~2, p.~6). It assumes a constant one-tick spread over the model's time scale (p.~6), and it does not apply to pro-rata markets, where queue position is irrelevant (n.~2, p.~2). Evidence: calibrated to Bank of America on Nasdaq over 22 trading days of August 2013 (Table~1, p.~19), and compared with backtests of artificial orders in nine large-tick stocks and ETFs (Table~2, August 2013; Table~3 reports averages over 30 trading days, p.~23).

In every symbol the front of the queue is worth more than the back; the model's gap between an order at the front (``touch value'') and one placed at the back (``order value'') is 0.22 ticks for BAC, 0.16 for CSCO, 0.17 for INTC, 0.11 for GE and F, 0.14 for EMM and 0.03 to 0.06 for PFE, PBR and EFA (Table~3; differences recomputed, \nolinkurl{dealer-queue-information-checks.py}). The authors call the gap for BAC and CSCO ``very large and comparable to the bid ask spread (${>}0.1$ ticks)'' (p.~24), and in the abstract ``of the same order of magnitude as the bid-ask spread''.

\qffield{Practice} The checkpoint of \nolinkurl{market-making-lob} asks for a queue-position fill model, and the overview's outcome asks how a replayed quoting policy depends on queue-position assumptions: an order replayed without its queue position overstates its fill probability and understates its adverse selection when it would in fact have joined the back of a long queue (Interpretation of the two components). The authors conclude that queue value ``rewards the trading speed of high-frequency firms'', disadvantaging traders without fast technology, and ask whether price--time priority suits large-tick assets (sec.~6, p.~25). Budish, Cramton and Shim note the same race for the top of the book when the tick is discrete (\hyperref[res:latency-arbitrage-arms-race]{latency-arbitrage arms race}); Harris's time precedence with a meaningful tick protects liquidity suppliers against quote matchers (\hyperref[def:order-anticipation]{order anticipation}).

\qffield{Pitfalls} The assertion was that speed buys queue priority and is most valuable ``when queues are thin, stale, or information-sensitive''. Priority and information sensitivity match the paper; ``stale'' matches its race after a price change. ``Thin'' does not: positional value is large where queues are long, in large-tick books, since in a thin queue every position is near the front (Interpretation of pp.~3--6). The paper is an author working paper; the catalogue records its 2016 version, while the copy read is the June 2017 revision, and a later journal version, if any, was not checked. Its analysis (sec.~3, the value recursion) was not checked in detail.

\qffield{See also} \hyperref[def:queue-imbalance]{queue imbalance}; \hyperref[mod:queue-reactive-intensities]{queue-reactive intensities}; \hyperref[mech:order-matching-pricing]{order matching and trade pricing} (time and pro-rata priority); \hyperref[def:tick-size]{tick size}; \hyperref[mod:make-take-fee-speed-race]{make/take fees and the speed race}.

\qffield{Further reading} Garriott and others on queuing (\nolinkurl{w3-garriott-queuing}) are an optional branch of \nolinkurl{market-making-lob}, unread.

\end{qfresult}

\begin{qfmodel}[Make/take fees and the speed race: high-frequency market makers and investors invest in speed, and the fee split that maximises trading (Foucault, Kadan and Kandel, as summarised by Lehalle and Laruelle)]\label{mod:make-take-fee-speed-race}
\qffield{Statement} Lehalle and Laruelle summarise a model of Foucault, Kadan and Kandel \citep[app.~A, sec.~A.2, pp.~250--253, page images of pp.~252--253]{alpha-laruelle}. There are $M$ high-frequency market makers and $N$ investors (market takers), modelled as buyers. An investor values the asset at $v_0+\Gamma$; a market maker sells at $v_0+s\Delta$, a spread of $s$ ticks of size $\Delta$; the venue charges $c_T$ to takers and $c_M$ to makers. Per trade the maker earns $\pi_M=s\Delta-c_M$, the investor $\pi_T=\Gamma-s\Delta-c_T$ and the venue $\bar c=c_M+c_T$. Market maker $m$ reaches the book at Poisson rate $\mu_m$ and investor $n$ at rate $\tau_n$; with the aggregate rates
\[ \bar\mu=\sum_{m=1}^{M}\mu_m,\qquad \bar\tau=\sum_{n=1}^{N}\tau_n,\qquad \mathcal R=\frac{\bar\mu\,\bar\tau}{\bar\mu+\bar\tau} \]
trades per unit time. Market makers race among themselves and investors among themselves, but not across sides. With quadratic costs of speed,
\[ \Pi_M(m)=\frac{\mu_m}{\bar\mu}\,\pi_M\,\mathcal R-\tfrac12\beta\mu_m^2,\qquad \Pi_T(n)=\frac{\tau_n}{\bar\tau}\,\pi_T\,\mathcal R-\tfrac12\gamma\tau_n^2,\qquad \Pi_E=\bar c\,\mathcal R, \]
and, once each participant has optimised its speed, the maker fee and taker fee satisfy
\[ c_M^\star=s\Delta-\frac{\Gamma-\bar c}{1+\bigl(\frac{M}{N}\cdot\frac{\gamma}{\beta}\bigr)^{1/3}},\qquad \bar c=c_T^\star+c_M^\star \]
(eq.~(A.4), p.~253). With few market makers the efficient maker fee can be negative (a maker rebate); with many market makers relative to investors, or when investors' speed is costlier ($\gamma$ large), the efficient taker fee can be negative (a taker rebate) (p.~253, Figure~A.2).

\qffield{Notation} $s\Delta$ the spread in currency; $\Gamma$ an investor's gain from trade; $\pi_M,\pi_T$ net gains per trade; $\mu_m,\tau_n$ arrival rates (per unit time), $1/\mu_m$ and $1/\tau_n$ times to market; $\bar\mu,\bar\tau$ \emph{sums} of rates; $\mathcal R$ the trade rate; $\beta,\gamma$ cost coefficients of speed; $\mu_m/\bar\mu$ and $\tau_n/\bar\tau$ shares of the races.

\qffield{Conditions} \emph{Derivation}. Taking $G=\Gamma-\bar c=\pi_M+\pi_T$ as fixed and letting each participant ignore its own effect on the aggregate rates (large $M$ and $N$), the first-order conditions give $\bar\mu(\bar\mu+\bar\tau)=M\pi_M\bar\tau/\beta$ and $\bar\tau(\bar\mu+\bar\tau)=N\pi_T\bar\mu/\gamma$; maximising $\mathcal R$ over the split $\pi_M$ then gives $\pi_M/\pi_T=(M\gamma/(N\beta))^{-1/3}$, that is $\pi_M=G/(1+(M\gamma/(N\beta))^{1/3})$, which is eq.~(A.4). Because $\bar c$ is fixed, the same split maximises the venue's revenue $\bar c\mathcal R$. Lehalle and Laruelle do not define ``efficient''; on this derivation it means trade-maximising for a given total fee, not welfare-maximising, since speed costs are not netted (Interpretation).

\emph{Against the original}. Foucault, Kadan and Kandel fix the platform's total fee $\bar c$ and choose the split to maximise $\Pi_e=(c_m+c_t)\mathcal R$ subject to $c_m+c_t=\bar c$, so the first-order condition equalises the marginal effects of the two fees on the trading rate (sec.~4, eqs.~(18)--(19), p.~23); in the large-market limit ($M,N\to\infty$ with $M/N=q$) Proposition~4 gives $c_m^*=\tfrac12\bigl(\Delta-2(L-\bar c)/(1+(qr)^{1/3})\bigr)$, $c_t^*=\bar c-c_m^*$ and $\pi_m^*=(L-\bar c)/(1+(qr)^{1/3})$, $\pi_t^*=(L-\bar c)/(1+(qr)^{-1/3})$, with $r=\gamma/\beta$ the takers' relative monitoring cost (eqs.~(22)--(24), p.~25, read as a page image) \citep[working paper, September 2009]{fkk-2009}.

That is Lehalle and Laruelle's (A.4) with their $s\Delta$ for the paper's $\Delta/2$ and $\Gamma$ for $L$, and it confirms both the large-population assumption and the reading of ``efficient'' above: the optimum is the trade-maximising (and so revenue-maximising) split for a given total fee. Welfare in the paper is the trading rate times the gains from trade less all monitoring costs, $W=\mathcal R L-MC_m-NC_t$, which rises as monitoring costs or fees fall (p.~20), a different objective. The make-take spread $c_m-c_t$ rises with the tick size, with $q$ and with $r$ (Corollary~6, p.~25).

Checked against exact symmetric equilibria with finite populations ($\beta=1$, $\gamma=2$, $G=1$): the trade-maximising maker margin is 0.4463 for $M=N=5$, 0.4429 for $M=N=50$ and 0.4425 for $M=N=5{,}000$, against 0.4425 from the formula (\nolinkurl{dealer-queue-information-checks.py}). \emph{Printed slips}: $\bar\tau$ is printed as $\sum_{m=1}^{M}\tau_m$ and the investor's cost as $\frac12\gamma\tau_m^2$, where the sum runs over the $N$ investors; the cost of speed is first written $\frac12\beta\mathbf T/t^2$ with $\mathbf T$ a trading day and then $\frac12\beta\mu_m^2$ in the gains (p.~252, page image).

The text's remark that ``$\Gamma-\bar c$ (the investor gain minus the fees) is obviously larger than $s\cdot\Delta$'' needs $\pi_T>c_M$, which the participation conditions $\pi_M,\pi_T>0$ do not imply ($s\Delta=1$, $c_M=0.9$, $c_T=0$, $\Gamma=1.05$ gives $\Gamma-\bar c=0.15$; checked).

\qffield{Practice} Lehalle and Laruelle introduce the model to address three questions raised by MiFID's market structure: the roles of high-frequency market makers, ``the possible negative externalities linked to their activity'', and ``the impact of the fee schedules on the efficiency of the PFP (Price Formation Process)'' (pp.~250--251). The model implies race shares and rebate cases (``maker rebates fit scarce or costly market-making; taker rebates fit crowded makers and costly investor access''), which match the text.

Their list of speed costs (co-location, feed handling, quoting engines and hedge automation for makers; broker connectivity, routing latency and algorithm responsiveness for investors), their claim that fast makers help price discovery when they update quotes and arbitrage stale venues but hurt it when speed ``mainly earns queue priority or fee rebates without adding durable liquidity'', and their diagnostic (segment by spread, maker count, investor urgency, queue stability, latency cost and whether rebates pay for fleeting quotes) are an interpretation: the model has no price discovery, no queues and no cancellations. That fast market makers keep prices consistent across venues is Lehalle and Laruelle's description (\hyperref[def:hft-strategies]{HFT strategies}); the fee comparison a router makes is \hyperref[res:fee-adjusted-venue-choice]{fee-adjusted venue choice}.

\qffield{Pitfalls} A simplified model describes $\bar\mu$ and $\bar\tau$ as averages; they are sums, which is what makes $\mu_m/\bar\mu$ a race share and $\mathcal R$ the trade rate. One definition gives eq.~(A.4) as ``the efficient maker fee'' without saying efficient for what; on the derivation above it is the split that maximises trading for a fixed total fee. Foucault, Kadan and Kandel's paper was read in its September 2009 working-paper version (above), not in the 2013 \emph{Journal of Finance} version; the pathway's empirical readings on fees were not read (\nolinkurl{ex-colliard-foucault-2012}, \nolinkurl{ex-malinova-park-2015}).

\qffield{See also} \hyperref[res:latency-arbitrage-arms-race]{latency-arbitrage arms race} (speed races that dissipate rents); \hyperref[res:queue-position-value]{queue-position value}; \hyperref[res:fee-adjusted-venue-choice]{fee-adjusted venue choice}; \hyperref[def:hft-strategies]{HFT strategies}; \hyperref[rule:us-equity-market-structure]{U.S.\ equity market structure} (access-fee caps).

\end{qfmodel}

\begin{qfmodel}[A toy model of the Flash Crash: an overflow of market orders depletes depth, a volume-following seller and a hot-potato game amplify it, and the time term of the solution (Lehalle and Laruelle)]\label{mod:depth-overflow-flash-crash}
\qffield{Statement} Lehalle and Laruelle model the mechanism they see in the Flash Crash: the usual trading rate at the bid or ask is near an equilibrium that depth renews; an agent consumed far more than depth could renew, emptied one side of the books, and kept its trading rate proportional to traded volume \citep[app.~A, sec.~A.5, pp.~261--262, page images of pp.~261--266]{alpha-laruelle}. With $F_t$ the arrival rate of sell market orders, traded volume $dV_t=F_t\,dt$, depth refilled at rate $\bar F$, overflow $dR_t=(F_t-\bar F)\,dt$, $r_t=F_t-\bar F$, and impact the inverse of depth, $I=1/D$,
\[ \frac{dI_t}{I_t}=\sigma_I\,dW_t+r_t\,dt \qquad\text{(A.12)},\qquad dP=\sigma^P dW^P+I^+dR^+-I^-dR^- \qquad\text{(A.13)} \]
(pp.~262--263). The \emph{hot-potato game}: high-frequency market makers who buy from a large seller hedge in another market against each other, so a fraction $\rho$ of traders re-trades, and a seller who follows market volume trades more; the rate becomes $r_t\simeq\frac{1-\rho^t}{1-\rho}\,r_0$ and the overflow $R_t\simeq\frac{t}{1-\rho}\,r_0$ (pp.~263--264). They then give
\[ I^-_t=I_0\exp\Bigl(\frac{t^2}{1-\rho}\,\frac{r_0}{2}-\tfrac12\sigma_I^2t+\sigma_IW^-_t\Bigr)\qquad\text{(A.14)} \]
and a price whose dominant part, $P_0-c\int\exp(t^2/u)\,dt$, ``is highly explosive, driving the price to zero faster than an exponential decay'' (p.~265). Their remedies: an obligation for market makers ``to refill the market depth quickly enough'' and a circuit breaker ``not only on price levels but also on market depth break events'' (pp.~265--266).

\qffield{Conditions} \emph{The time term}. (i) In $r_t=\frac{1-\rho^t}{1-\rho}r_0=(1+\rho+\dots+\rho^{t-1})r_0$, $t$ counts rounds of re-trading: it gives $r_0$ after one round and 0, not $r_0$, at $t=0$. In clock time on $[0,1]$, the time axis of Figures~A.4--A.5, $\rho^t\ge\rho$ and $r_t\le r_0$: no amplification. Using it in continuous time needs a round length $\delta$, with $t/\delta$ rounds by time $t$, which the text does not give. (ii) The overflow after $k$ rounds is $\frac{r_0}{1-\rho}\bigl(k-\frac{\rho(1-\rho^k)}{1-\rho}\bigr)$; the text's $kr_0/(1-\rho)$ is its large-$k$ form and overstates it by up to $\rho r_0/(1-\rho)^2$ (3.75$\,r_0$ at $\rho=0.6$). (iii) The exponent $\frac{t^2}{1-\rho}\frac{r_0}{2}$ of (A.14) equals $\int_0^tR_s\,ds$ with $R_s=s\,r_0/(1-\rho)$: (A.14) solves $dI/I=\sigma_IdW+R_t\,dt$, with the overflow, not (A.12) with the rate, whose solution with $r_t=r_0/(1-\rho)$ has the linear exponent $r_0t/(1-\rho)$.

The text's gloss supports the overflow reading: ``when more volume is traded than expected, R is large, thus the impact I is large'' (p.~262). Checked by simulation: $E\log(I_1/I_0)$ is 2.459 against 2.455 with the rate and 1.210 against 1.205 with the overflow ($\rho=0.6$, $r_0=1$, $\sigma_I=0.3$, \nolinkurl{dealer-queue-information-checks.py}). (iv) $r\,dt$ is a quantity while $dI/I$ is a number, so a scale (per unit quantity for $r\,dt$, per unit quantity and time for $R\,dt$) is implicit. (v) The amplification converges only for $\rho<1$.

\qffield{Practice} The \emph{liquidity spiral} (a large order consumes depth, widens spreads and raises volatility; providers cut size or cancel; the next order trades through less depth) is the verbal mechanism of this model; the term and its state variable, ``confidence in replenishment'', are the interpretation's, not the source's. The illustrative control list (rate controls on volume-following execution, depth controls, price controls, cross-market coordination and real-time detection of overflow and broken references) combines Lehalle and Laruelle's two remedies with the lessons of the 2010 staff report (\hyperref[ev:may-2010-flash-crash]{the May 2010 Flash Crash}): pauses such as the CME Stop Logic halt and single-stock circuit breakers (\hyperref[mech:circuit-breakers-trading-halts]{circuit breakers}). It is a synthesis, not a list from one source.

\qffield{Pitfalls} The simplified model block reproduces (A.12) faithfully; its statement that positive overflow makes inverse depth ``compound upward'' holds for (A.12), but the super-exponential fall of the text needs the overflow reading of (iii). The \emph{hot-potato amplification} says that ``a fraction $\rho$ of traders follows volume mechanically''; in the text $\rho$ is ``the proportion of traders playing this hot potato game'' (p.~264), and the volume follower is the seller. The model is illustrative: its parameters are not estimated, its figures plot depth with a sign convention (bid depth negative) not stated in the text, and the staff report describes the hot potato as high-frequency traders re-trading E-mini contracts among themselves, with cross-market arbitrageurs carrying the pressure to equities, not as market makers hedging across markets with each other. The staff report was read only in its executive summary, and the pathway's Flash Crash readings (\nolinkurl{w5-kirilenko-flash-crash}, \nolinkurl{w5-menkveld-yueshen-flash}) not at all.

\qffield{See also} \hyperref[ev:may-2010-flash-crash]{the May 2010 Flash Crash}; \hyperref[mech:bubble-crash-dynamics]{bubbles and crashes}; \hyperref[def:liquidity-dimensions]{liquidity and market quality} (resiliency); \hyperref[mech:circuit-breakers-trading-halts]{circuit breakers and trading halts}; \hyperref[mod:propagator]{propagator model}.

\end{qfmodel}

\begin{qfevidence}[The Flash Crash of 6 May 2010: a volume-targeting sell algorithm, a hot-potato volume effect in E-mini futures, withdrawal of liquidity in equities and trades at stub quotes (CFTC and SEC staff report, 2010)]\label{ev:may-2010-flash-crash}
\qffield{Statement} The CFTC and SEC staffs' findings, reported on 30 September 2010 \citep[executive summary, pp.~1--7]{sec-cftc-may6-2010}; the Commissions ``expressed no view'' on them (cover):
\begin{itemize}\setlength{\itemsep}{0pt}
\item \emph{Setting.} An unusually turbulent day; by 2:30 p.m.\ buy-side depth in the E-mini S\&P~500 futures had fallen from nearly \$6 billion to \$2.65 billion and in SPY from about \$275 million to \$220 million (pp.~1--2).
\item \emph{Trigger.} At 2:32 p.m.\ a large fundamental trader, a mutual fund complex, began selling 75{,}000 E-mini contracts (about \$4.1 billion) as a hedge, through an algorithm targeting 9 percent of the previous minute's volume ``but without regard to price or time''; a previous program of that size had taken more than five hours, this one about 20 minutes (p.~2).
\item \emph{Absorption and the hot potato.} High-frequency traders and other intermediaries, fundamental buyers and cross-market arbitrageurs, who sold SPY or stocks against the E-mini, absorbed the selling; HFTs, net long about 3{,}300 contracts, sold about 2{,}000 between 2:41 and 2:44 p.m.\ while trading nearly 140{,}000; the algorithm sped up as volume rose; then, between 2:45:13 and 2:45:27, HFTs traded over 27{,}000 contracts, about 49 percent of volume, while buying only about 200 net, ``a `hot-potato' volume effect'' (p.~3). Buy-side E-mini depth fell to about \$58 million, under 1 percent of the morning's, and the E-mini fell more than 5 percent and SPY more than 6 percent between 2:41 and 2:45:27 (pp.~3--4).
\item \emph{Pause and recovery.} CME's Stop Logic paused the E-mini for five seconds at 2:45:28; prices stabilised and recovered (p.~4).
\item \emph{Equities.} Liquidity providers' automated systems paused; some widened quotes, others cut size or withdrew, weighing data integrity, risk and position limits, P\&L and the risk of broken trades; internalisers routed retail orders to exchanges; stop-loss orders met reduced buying interest; with liquidity gone, market orders executed against stub quotes at prices ``as low as one penny or as high as \$100,000''; over 20{,}000 trades in more than 300 securities at prices 60 percent or more away from their 2:40 p.m.\ prices were broken; by about 3:00 p.m.\ most securities traded near consensus values again (pp.~4--6).
\end{itemize}

\qffield{Conditions} Lessons drawn by the staff: under stress, the automated execution of a large sell order ``can trigger extreme price movements, especially if the automated execution algorithm does not take prices into account''; ``high trading volume is not necessarily a reliable indicator of market liquidity''; derivatives and securities markets are interconnected through index products; one participant's withdrawal may not matter, but ``a liquidity crisis can develop if many market participants withdraw at the same time''; and pausing a market can let participants reassess (pp.~3 and~6). NYSE's liquidity replenishment points, Nasdaq's self-help declaration against NYSE Arca and delays in NYSE data were examined, and ``none of these factors played a dominant role'' (n.~10, p.~5).

The arithmetic is consistent: \$4.1 billion for 75{,}000 contracts and \$1.9 billion for 35{,}000 imply an index level near 1{,}090 at 50 dollars a point; 27{,}000 contracts at 49 percent imply about 55{,}000 traded in 14 seconds, 135 times the HFTs' net purchase (\nolinkurl{dealer-queue-information-checks.py}).

\qffield{Practice} Lehalle and Laruelle's account \citep[chap.~2, sec.~2.5.2, pp.~177--183]{alpha-laruelle} agrees on the trigger and adds Procter \& Gamble's fall from \$61.50 to \$39.37 between 14:40 and 14:47 and its volume-weighted spread above \$2 (pp.~179--180). It describes the hot potato as high-frequency counterparts hedging in equities and then in futures in a loop, gives NYSE's liquidity replenishment points a possible aggravating role (pp.~180--181) and concludes that ``Every layer of the market contributed to the crash'' (p.~183). The staff report places the hot potato within the E-mini and found the replenishment points not dominant. The case study (a large automated futures sell program met thin liquidity; HFTs recycled inventory; aggressive flow moved into cash equities; spreads widened; routers or venues paused) matches the report.

Of the \emph{shock propagation channels}, cross-asset hedging, routing changes, quote withdrawal and correlated constraints (risk and position limits, pauses) are in the report; \emph{reference-price dependency} is documented by Lehalle and Laruelle for the Euronext outage of 13 October 2010, when midpoint trading stopped and orders pegged to the primary's quotes were cancelled or paused (p.~184; \hyperref[ev:european-fragmentation-evidence]{European fragmentation evidence}), not in the report's summary.

\emph{Measured HFT behaviour}. Kirilenko, Kyle, Samadi and Tuzun use CFTC audit-trail transactions in the June 2010 E-mini for 3--6 May 2010; of over 15,000 accounts, intraday intermediaries are accounts that traded at least ten contracts on one of 3--5 May and whose end-of-day positions averaged no more than 5 percent of daily volume, with minute-by-minute deviations from them of no more than 0.5 percent, over 3--5 May, the 16 most active of them being ``High Frequency Traders'' and the rest ``Market Makers'', with categories fixed from 3--5 May into 6 May (PDF pp.~5 and 11--14) \citep[accepted manuscript]{w5-kirilenko-flash-crash}. High Frequency Traders had 34.22 percent of volume on 3--5 May and 28.57 percent on 6 May, Market Makers 10.49 and 9.00 (Table~II, PDF p.~35).

The intermediaries' combined net inventories never exceeded 6,000 contracts over the four days, an order of magnitude below the 75,000-contract sale, and both groups accumulated net long positions as prices fell; High Frequency Traders' second-by-second inventory changes co-move \emph{positively} with contemporaneous and lagged price changes, a relationship that ``did not change when prices fell during the Flash Crash'', whereas Market Makers' inventory changes are negatively related to price changes and their relationship did change (PDF pp.~5--6). So in this study HFTs neither caused the crash nor changed behaviour during it; the hot-potato trading among them is the staff report's.

\emph{The report's detailed sections} (sec.~I, printed pp.~12--17, read 2026-10-16) add: E-mini buy-side resting orders fell below 1{,}050 contracts by 2:45:28, while SPY kept about 600{,}000 shares (about 1{,}200 contract equivalents) of buy depth, and near-inside depth was not affected; the Stop Logic pause is triggered to keep stop-loss cascades inside a ``no bust'' range of six index points, about 0.6 percent on the day (n.~20, p.~12).

The trader categories behind the summary, classified on 3--5 May, include HFTs as the top 3 percent of intermediaries by number of trades, 16 of over 15{,}000 accounts trading 1{,}455{,}000 contracts on 6 May, almost a third of volume, while rarely holding more than 3{,}000 contracts either way; the section is ``based in part'' on Kirilenko et al.\ (n.~22, pp.~13 and~15). Of the 35{,}000 contracts the sell algorithm sold from 2:32 to 2:45, about 18{,}000 were executed aggressively and 17{,}000 passively (n.~27), and of the remaining 40{,}000, sold from 2:45 to 2:51 as prices rose, about 24{,}000 aggressively (n.~28); over 2:32--2:45 all fundamental sellers sold more than 80{,}000 contracts net against about 50{,}000 bought by fundamental buyers, a 30{,}000 imbalance about 15 times that of the same interval on the previous three days; in the rebound HFTs ``did not significantly alter their trading strategy'' while nearly half the intermediaries withdrew (pp.~15--16).

Fifteen cross-market firms interviewed, net buyers of more than 100{,}000 E-minis between 2:00 and 3:00 p.m., saw the E-mini as the primary price-discovery product, bought it against SPY and stock baskets during the fall, and many had stopped their strategies by 2:45 p.m. (sec.~I.5, pp.~16--17). Menkveld and Yueshen's account of the broken E-mini--SPY arbitrage link is available in the SSRN working-paper version (Pitfalls).

\qffield{Pitfalls} These are the staffs' findings from their reconstruction; the executive summary and the E-mini and cross-market sections (sec.~I) were read, the equities sections (secs.~II--III) through the summary and, on 2026-10-16, their parts on internalisers and stub quotes. Those confirm the summary's equities bullet: OTC market makers ``paused or halted internalization'', some only on sell orders as position limits were hit, citing sell pressure, data integrity, latencies and P\&L limits, and routed orders to the exchanges instead \citep[sec.~II.2.e, printed p.~58]{sec-cftc-may6-2010}; many retail market and stop-loss orders were converted by internalisers into limit orders that followed the price down to stub-level bids; and the ``over 20{,}000 trades representing 5.5 million shares'' executed more than 60 percent from their 2:40 p.m.\ value were broken under the clearly erroneous rules, almost two-thirds of the cancelled shares having traded below \$1.00 and about 5 percent above \$100 (sec.~II.3.a, printed p.~65).

Later accounts other than the two papers named here were not read. \emph{Cross-market arbitrage} (read 2026-10-16 in the SSRN version dated 15 November 2017, earlier titled ``Anatomy of the Flash Crash''; the 2019 \emph{Management Science} version was not compared). Menkveld and Yueshen proxy cross-arbitrage by the dispersion of bids and of asks across venues, using Nanex quote data and Waddell \& Reed's own trade records, and find that arbitrage between the E-mini and SPY ``weakened at 14:15 and broke entirely at 14:44, one minute before the halt'' \citep[SSRN version, printed pp.~2--3 and~5--7]{w5-menkveld-yueshen-flash}; in the 50 most crashed S\&P~500 stocks, those whose arbitrage broke first crashed first, with controls for volatility and order imbalance, but ``We however are unable to identify the root cause for these breakdowns themselves'' (pp.~3 and~15).

The large seller traded most of its 75,000 contracts after the breakdown, and at the widest the E-mini ask was more than 100 basis points below the SPY bid (p.~3). A Grossman--Miller calibration, in which a disconnected venue leaves the seller facing half the market makers, puts its overpayment for liquidity at \$98.6 million to \$229.8 million, two to four times its quarterly operating income (pp.~3--4): a model-based estimate, not an accounting figure. Their online appendix dates the break points that separate cointegrated from non-cointegrated E-mini and SPY prices, so a VECM is fitted before the break and a VAR after \citep[secs.~A--B]{menkveld-yueshen-oa}. These are associations and a calibration, read alongside the staff's reconstruction, not a causal identification.

The sell algorithm was about half passive: ``taking liquidity'' describes its effect on net demand, not every fill. The Kirilenko et al.\ figures are from the accepted manuscript in Cambridge's repository, not the published pages. The rules changed afterwards (single-stock circuit breakers from June 2010, limit up--limit down, revised market-wide circuit breakers, the stub-quote ban: \hyperref[rule:us-equity-market-structure]{U.S.\ equity market structure}). ``Hot potato'' names a volume effect among intermediaries, not a measure of HFT's net selling.

\qffield{See also} \hyperref[mod:depth-overflow-flash-crash]{toy model of the Flash Crash}; \hyperref[mech:circuit-breakers-trading-halts]{circuit breakers}; \hyperref[def:hft-strategies]{HFT strategies}; \hyperref[ev:index-futures-price-discovery]{index futures and price discovery}; \hyperref[mech:venue-competition-fragmentation]{venue competition and fragmentation}.

\end{qfevidence}
